\subsection{Collection of Existing Results on Stochastic TSP}\label{sec:results-collection}

This section collects all known results on stochastic TSP variants (a posteriori, adaptive, a priori) from the literature.
We use the notation from Section~2 of this thesis:
$\optpost = \E_A[\opt(A)]$ (a posteriori),
$\optadapt$ (optimal adaptive policy),
$\optprior$ (optimal a priori master tour).
The fundamental inequality is $\optpost \le \optadapt \le \optprior$.
We define:
\begin{itemize}
    \item \textbf{Obliviousness gap}: $\sup_{\II} \f{\optprior}{\optpost}$ \quad (a priori vs.\ a posteriori)
    \item \textbf{Adaptivity gap}: $\sup_{\II} \f{\optprior}{\optadapt}$ \quad (a priori vs.\ adaptive)
    \item \textbf{Clairvoyance gap}: $\sup_{\II} \f{\optadapt}{\optpost}$ \quad (adaptive vs.\ a posteriori)
\end{itemize}
Note: The literature on a priori TSP (e.g., \cite{Christalla2025}) uses ``adaptivity gap'' to mean what we call the obliviousness gap. We follow the terminology of our thesis (Section~3).

%% ============================================================
\subsubsection{Symmetric Metrics}
%% ============================================================

\paragraph{Gap Bounds.}

\begin{center}
\renewcommand{\arraystretch}{1.4}
\resizebox{\textwidth}{!}{%
\begin{tabular}{l c c l}
    \hline
    \textbf{Gap} & \textbf{Lower bound} & \textbf{Upper bound} & \textbf{Source} \\
    \hline
    Obliviousness $\f{\optprior}{\optpost}$
        & ---
        & $\le 3$
        & Shmoys--Talwar~\cite{Shmoys2008} \\[3pt]
    Adaptivity $\f{\optprior}{\optadapt}$
        & ---
        & $\le 3$ (inherited)
        & Since $\optadapt \ge \optpost$ \\[3pt]
    Clairvoyance $\f{\optadapt}{\optpost}$
        & ---
        & $\le 3$ (inherited)
        & Since $\optprior \ge \optadapt$ \\
    \hline
\end{tabular}%
}
\end{center}

\begin{remark}
    The obliviousness gap bound $\le 3$ from \cite{Shmoys2008} is the best known.
    It follows from their sampling algorithm: draw a random active set~$S$,
    build an MST on~$S$,
    connect each non-sampled vertex to its nearest neighbor in~$S$,
    and double the resulting tree.
    The expected cost is bounded by
    $2(\E[\text{MST}(S)] + \E[\sum_j D_j(S)])
    \le 2(\optpost + \optpost)$.
    Since for any master tour~$T$,
    \[
        \E[c(T|_A)] \le \E[\text{cost of sample-based tour}],
    \]
    the ratio $\optprior / \optpost \le 3$ follows from a
    more refined analysis using LP lower bounds.
    No non-trivial lower bounds on the obliviousness, adaptivity,
    or clairvoyance gaps are known for symmetric metrics
    (i.e., it is open whether any of these gaps can be $> 1 + \eps$
    for some constant $\eps$).
\end{remark}

On tree metrics, the situation is completely resolved:

\begin{center}
\renewcommand{\arraystretch}{1.4}
\begin{tabular}{l c l}
    \hline
    \textbf{Gap} & \textbf{Value} & \textbf{Source} \\
    \hline
    Obliviousness $\f{\optprior}{\optpost}$ on trees & $= 1$ & Schalekamp--Shmoys~\cite{SchalekampShmoys2008} \\
    \hline
\end{tabular}
\end{center}

\begin{remark}
    On tree metrics, any shortcutting of the Euler tour traversal of the tree gives an optimal tour for \emph{every} subset $S$ simultaneously (Theorem~2 in \cite{SchalekampShmoys2008}). Thus the universal TSP ratio is exactly~1, implying all three gaps equal~1 on tree metrics.
\end{remark}

\paragraph{Approximation Algorithms for A Priori TSP (Symmetric).}

These algorithms compute a master tour $T$ and bound $\E[c(T|_A)]$ relative to $\optprior$.

\begin{center}
\renewcommand{\arraystretch}{1.4}
\resizebox{\textwidth}{!}{%
\begin{tabular}{l c c l}
    \hline
    \textbf{Result} & \textbf{Ratio vs.\ $\optprior$} & \textbf{Type} & \textbf{Source} \\
    \hline
    Sampling + MST doubling & $4$ & Randomized & Shmoys--Talwar~\cite{Shmoys2008}, Thm~10 \\
    CFL + derandomization & $8$ & Deterministic & Shmoys--Talwar~\cite{Shmoys2008}, Thm~13 \\
    Sampling, $\sigma\!=\!0.663$, $\al\!=\!1.5$ & $< 3.1$ & Randomized & Blauth et al.~\cite{Blauth2024}, Cor~4 \\
    Master route + Karlin--Klein--OG & $5.9$ & Deterministic & Blauth et al.~\cite{Blauth2024}, Cor~7 \\
    Tree embedding & $O(\log n)$ & Rand., dist.-free & Schalekamp--Shmoys~\cite{SchalekampShmoys2008}, Cor~4 \\
    \hline
\end{tabular}%
}
\end{center}

\begin{remark}[Details on Shmoys--Talwar~\cite{Shmoys2008}]
    The $4$-approximation works as follows: sample a random set $S$ from the activation distribution conditioned on $|A| \ge 2$; build MST on $S$; for each $j \notin S$, connect $j$ to its nearest neighbor in $S$; double the resulting tree. The key insight is Lemma~7: $\E_A[D_j(A)] \le \E_A[D_j(A) \mid j \in A]$, where $D_j(A)$ is the nearest-neighbor distance from $j$ in $A$. This uses Facts~3 and~4: $\optpost \ge \E[\text{MST}(A)]$ and $\optpost \ge \E[\sum_{j \in A} D_j(A)]$.

    The $8$-approximation uses the same structure but derandomizes via a Connected Facility Location LP relaxation. The LP optimum is $\le \f{3}{2}\optprior$ (Lemma~11), and conditional expectations give a deterministic set $S$ with $2 \cdot c_{\mathrm{ST}}(S) + c_R(S) \le 4\optprior$; doubling gives $8\optprior$.
\end{remark}

\begin{remark}[Details on Blauth et al.~\cite{Blauth2024}]
    Their approach uses a \emph{master route solution}: select a nonempty set $S$, compute a TSP tour on $S$, and connect each vertex in $V \setminus S$ to its nearest point in $S$ by a pair of parallel edges. The key results are:

    \begin{itemize}
        \item \textbf{Master route ratio} (Theorem~6): The supremum of $\min_S \text{MR}(S) / \optprior$ over all instances with depot is at least $1/(1 - e^{-1/2}) > 2.541$ and at most $2.584$. They conjecture the exact value is $1/(1 - e^{-1/2})$.

        \item \textbf{Sampling lower bound} (Theorem~2): No matter how the sampling function $f$ is chosen, the sampling algorithm cannot achieve ratio better than $2.655$ (even with optimal TSP subroutine) or $3.049$ (with a $1.4999$-approximate TSP subroutine).

        \item \textbf{Depot reduction} (Theorem~1): If there is a $\rho$-approximation for a priori TSP with depot (for $\rho \ge 3$), then there is a $(\rho + \eps)$-approximation for general instances. Loss is arbitrarily small.

        \item \textbf{Deterministic algorithm} (Corollary~7): Combines the master route ratio $< 2.6$ with the deterministic $3/2$-approximation for TSP from Karlin--Klein--Oveis Gharan and van Zuylen's framework (Theorem~35: ratio $\le 2 + \al \cdot \rho$ where $\al = 3/2$, $\rho < 2.6$, giving $2 + 3.9 = 5.9$).
    \end{itemize}
\end{remark}

\begin{remark}[Schalekamp--Shmoys {\cite{SchalekampShmoys2008}}]
    Their $O(\log n)$ result is distribution-free: the same tour works
    for \emph{any} distribution over active sets.
    The algorithm uses FRT stochastic tree embeddings
    with expected distortion~$O(\log n)$:
    sample a random tree metric, compute the optimal tree
    traversal (optimal for every subset by Theorem~2),
    and shortcut back to the original metric.
    The $O(\log n)$ bound cannot be improved using tree embedding
    techniques (by a result of Leighton--Rao).
\end{remark}

\paragraph{Approximation Algorithms for Adaptive TSP (Symmetric).}

These algorithms compute an adaptive policy and bound its expected cost relative to the optimal adaptive policy.

\begin{center}
\renewcommand{\arraystretch}{1.4}
\begin{tabular}{l c c l}
    \hline
    \textbf{Result} & \textbf{Ratio vs.\ $\optadapt$} & \textbf{Model} & \textbf{Source} \\
    \hline
    Scenario-halving + GSO & $O(\log^2 n \cdot \log m)$ & Visit-to-reveal & GNR~\cite{GuptaNagarajanRavi2017}, Thm~2 \\
    \hline
\end{tabular}
\end{center}

\begin{remark}[Details on Gupta--Nagarajan--Ravi~\cite{GuptaNagarajanRavi2017}]
    This paper studies a \emph{different} adaptive TSP model: the activation status of vertex $v$ is revealed only upon \emph{visiting} $v$ (not upon calling). This is a harder model than ours (the salesperson has strictly less information per unit of cost). The distribution is given explicitly with support size $m$.

    The algorithm has the following structure:
    \begin{enumerate}
        \item Reduce Adaptive TSP to the \emph{Isolation Problem}: given $m$ scenarios, find a minimum-cost adaptive strategy that identifies which scenario is true.
        \item Reduce Isolation to \emph{Group Steiner Orienteering (GSO)}: find a walk of bounded length that hits the maximum number of groups.
        \item Solve GSO via LP relaxation and tree embeddings: $(4, O(\log^2 n))$-bicriteria approximation (Theorem~18).
        \item The scenario-halving argument reduces live scenarios by factor $7/8$ per round, so $O(\log m)$ rounds suffice.
    \end{enumerate}

    The Isolation-to-AdapTSP reduction (Lemma~17) loses an additive $3/2$: use the $\al$-approximate isolation tree to identify the realized scenario, then traverse a $3/2$-approximate TSP tour on the active set, total cost $(\al + 3/2) \cdot \optadapt$.

    \textbf{Hardness}: $\Omega(\log^{2-\eps} n)$ on tree metrics, by reduction from Group Steiner Tree (Theorem~27).

    This paper does not bound any gap between adaptive and a priori/a posteriori costs. It approximates the optimal adaptive policy directly.
\end{remark}

%% ============================================================
\subsubsection{Asymmetric Metrics}
%% ============================================================

\paragraph{Gap Bounds.}

\begin{center}
\renewcommand{\arraystretch}{1.4}
\resizebox{\textwidth}{!}{%
\begin{tabular}{l c c l}
    \hline
    \textbf{Gap} & \textbf{Lower bound} & \textbf{Upper bound} & \textbf{Source} \\
    \hline
    Obliviousness $\f{\optprior}{\optpost}$
        & $\Omega\!\left(\f{n^{1/4}}{\log n}\right)$
        & $O(\sqrt{n})$
        & CPT~\cite{Christalla2025}, Thms~1.1,~1.2 \\[6pt]
    Adaptivity $\f{\optprior}{\optadapt}$
        & $\Omega\!\left(\f{n^{1/4}}{\log n}\right)$
        & $O(\sqrt{n})$
        & Same LB instance; UB inherited \\[6pt]
    Clairvoyance $\f{\optadapt}{\optpost}$
        & ---
        & $O(\sqrt{n})$ (inherited)
        & \textbf{Open}: is this $O(1)$? \\
    \hline
\end{tabular}%
}
\end{center}

\begin{remark}[Details on the $\Omega(n^{1/4} / \log n)$ lower bound~\cite{Christalla2025}]
    The lower bound uses \emph{asymmetric grid instances} $G_k$: vertices are placed in a $k \times k$ grid, with $k^2$ vertices and activation probability $p(v) = 1/k$ for all $v$. The asymmetric distances are designed so that:
    \begin{itemize}
        \item The optimal a posteriori cost satisfies
        $\E[\text{TSP}(A, c_k)] \le 5k \log k$
        (Lemma~3.6, using Chernoff bounds on the maximum column occupancy;
        Lemma~3.5 shows
        $\text{TSP}(A, c_k) = k \cdot \max_j |A \cap C_j|$).
        \item The optimal a priori tour satisfies (Lemma~3.4):
        \[
            \optprior(V_k, c_k, p_k) \ge \f{k^{3/2}}{2^9 \cdot e^4}.
        \]
    \end{itemize}
    Thus for $n = k^2$ vertices, the obliviousness gap is at least
    \[
        \f{k^{3/2}/(2^9 e^4)}{5k \log k} = \Omega\!\left(\f{\sqrt{k}}{\log k}\right) = \Omega\!\left(\f{n^{1/4}}{\log n}\right).
    \]
    As noted in our thesis (Section~3), this construction also implies the adaptivity gap $\f{\optprior}{\optadapt} = \Omega(n^{1/4} / \log n)$ because $\optadapt \le \optprior$.
\end{remark}

\begin{remark}[Details on the $O(\sqrt{n})$ upper bound~\cite{Christalla2025}]
    Theorem~1.1 gives a deterministic polynomial-time algorithm computing a tour $T$ with
    \[
        \E[c(T|_A)] \le O(\sqrt{n}) \cdot \E[\text{TSP}(A, c)].
    \]
    The algorithm splits vertices into high-probability
    ($p(v) \ge 1/\sqrt{n}$) and low-probability
    ($p(v) < 1/\sqrt{n}$) sets.
    For high-probability vertices, it uses an $\al$-approximate ATSP tour
    costing
    $\al \cdot \cei{1/p_{\min}} \cdot \optpost \le \al\sqrt{n} \cdot \optpost$
    (Proposition~4.7).
    For low-probability vertices, any Hamiltonian cycle has expected cost
    $\le 6 \cdot p(V) \cdot \optpost$ (Proposition~4.8).
    Since $p(V_2) \le \sqrt{n}$, the combined cost is
    $O(\sqrt{n}) \cdot \optpost$.
\end{remark}

\paragraph{Approximation Algorithms for A Priori TSP (Asymmetric).}

\begin{center}
\renewcommand{\arraystretch}{1.4}
\begin{tabular}{l c c l}
    \hline
    \textbf{Result} & \textbf{Ratio} & \textbf{Compared to} & \textbf{Source} \\
    \hline
    High/low probability split & $O(\sqrt{n})$ & $\optpost$ & CPT~\cite{Christalla2025}, Thm~1.1 \\
    Hop-ATSP reduction & $O(\log^8 n)$ & $\optprior$ & CPT~\cite{Christalla2025}, Thm~1.3 \\
    \hline
\end{tabular}
\end{center}

\begin{remark}[Details on the $O(\log^8 n)$ quasi-polynomial-time algorithm~\cite{Christalla2025}]
    Theorem~1.3 gives a randomized algorithm with running time
    $n^{O(\log n)}$ that computes a tour~$T$ with
    $\E[c(T|_A)] \le O(\log^8 n) \cdot \optprior$.
    This is remarkable because it achieves an approximation ratio
    \emph{below} the obliviousness gap (which is $\Omega(n^{1/4}/\log n)$).
    The algorithm works through a chain of reductions:
    \begin{enumerate}
        \item Reduce a priori TSP to \emph{Hop-ATSP} via uniform
        activation probabilities (Theorem~2.2):
        an $\al(n)$-approx\-imation for Hop-ATSP yields an
        $O(\al(5n^2))$-approx\-imation for a priori TSP.
        \item Reduce Hop-ATSP to hierarchically ordered
        instances via directed low-diameter
        decompositions (Thm~2.3; loss
        $O(\log n \cdot \log\log n)$,
        depth $L = O(\log n)$).
        \item Reduce hierarchically ordered Hop-ATSP to a \emph{Path Covering}
        problem (Theorem~6.1): loss $O(\log^6 n) \cdot \al(n)$.
        \item Solve Path Covering in quasi-polynomial time via dynamic programming (Theorem~7.1): $O(\log^2 n)$-approximation in time $n^{O(\log n)}$.
    \end{enumerate}
    Combining: $O(\log^2 n) \cdot O(\log^6 n) = O(\log^8 n)$. The exponent $8$ is not optimized.

    A key technical result (Theorem~4.16) shows that for uniform activation
    probability~$\delta$, the expected a priori cost $\E[c(T|_A)]$ and
    $\delta^2 \cdot c^{(k)}(T)$ (the $k$-hop cost with $k = \flr{1/\delta}$)
    differ by at most a constant factor.
\end{remark}

%% ============================================================
\subsubsection{Open Problems and Gaps in Knowledge}
%% ============================================================

The main open questions for stochastic TSP are:

\begin{enumerate}
    \item \textbf{Clairvoyance gap for asymmetric metrics}: Is $\f{\optadapt}{\optpost} = O(1)$? Or does it grow with $n$? This is the central question of this thesis.

    \item \textbf{Tightening asymmetric obliviousness gap}: Currently $\Omega(n^{1/4}/\log n) \le \f{\optprior}{\optpost} \le O(\sqrt{n})$. The exact growth rate is unknown.

    \item \textbf{Symmetric gap lower bounds}: No non-trivial lower bounds are known for any of the three gaps on symmetric metrics. Is there an instance family where $\f{\optprior}{\optpost} > 1 + \eps$ for some constant $\eps > 0$?

    \item \textbf{Polynomial-time algorithm below the obliviousness gap (asymmetric)}: The $O(\log^8 n)$-approximation of~\cite{Christalla2025} runs in quasi-polynomial time $n^{O(\log n)}$. Can this be made polynomial?
\end{enumerate}

%% ============================================================
\subsubsection{Results on Related Stochastic Optimization Problems}
%% ============================================================

The following results on related problems provide context, techniques, and analogies for stochastic TSP.

\paragraph{Stochastic $k$-TSP.}

In Stochastic $k$-TSP, each vertex has a random reward; the salesperson must collect total reward $\ge k$ while minimizing expected travel cost.

\begin{center}
\renewcommand{\arraystretch}{1.4}
\begin{tabular}{l c l}
    \hline
    \textbf{Result} & \textbf{Bound} & \textbf{Source} \\
    \hline
    Adaptivity gap lower bound & $\ge e$ & Ene--Nagarajan--Saket~\cite{EneNagarajanSaket2017} \\
    Adaptivity gap upper bound & $O(1)$ & Jiang--Li--Liu--Singla~\cite{JiangLiLiuSingla2020} \\
    Adaptive $O(\log k)$-approx & $O(\log k)$ vs.\ $\optadapt$ & Ene--Nagarajan--Saket~\cite{EneNagarajanSaket2017} \\
    Nonadaptive $O(1)$-approx & $O(1)$ vs.\ $\optadapt$ & Jiang--Li--Liu--Singla~\cite{JiangLiLiuSingla2020} \\
    \hline
\end{tabular}
\end{center}

\begin{remark}
    The adaptivity gap for Stochastic $k$-TSP is $O(1)$, meaning adaptivity provides at most a constant-factor benefit. This contrasts with stochastic TSP on asymmetric metrics where the obliviousness gap grows with $n$.
\end{remark}

\paragraph{Stochastic Probing.}

In the stochastic probing framework~\cite{GuptaNagarajan2013}, elements have random activation; probing reveals status. Solutions must satisfy inner (selected) and outer (probed) constraints.

\begin{center}
\renewcommand{\arraystretch}{1.4}
\resizebox{\textwidth}{!}{%
\begin{tabular}{l c l}
    \hline
    \textbf{Constraint type} & \textbf{Adaptivity gap} & \textbf{Source} \\
    \hline
    Matroid inner + matroid outer, submod.\ obj. & $O(1)$ & Gupta--Nagarajan--Singla~\cite{GuptaNagarajanSingla2017} \\
    Matroid inner + prefix-closed outer & $O(\log n)$ & Gupta--Nagarajan--Singla~\cite{GuptaNagarajanSingla2016} \\
    Single matroid inner, submod.\ obj. & $= 2$ & Bradac--Singla--Zuzic~\cite{BradacSinglaZuzic2019} \\
    \hline
\end{tabular}%
}
\end{center}

\paragraph{Stochastic Knapsack.}

In stochastic knapsack, items have random sizes; the knapsack has capacity~1. The adaptivity gap measures the benefit of choosing items dynamically vs.\ committing to a fixed order.

\begin{center}
\renewcommand{\arraystretch}{1.4}
\begin{tabular}{l c c l}
    \hline
    \textbf{Variant} & \textbf{Lower bound} & \textbf{Upper bound} & \textbf{Source} \\
    \hline
    Non-Risky SK & $1.37$ & $4$ & DGV~\cite{DeanGoemansVondrak2008}, BTC~\cite{BarakTalgamCohen2026} \\
    Risky SK (All-or-None) & $2$ & $8.47$ & BTC~\cite{BarakTalgamCohen2026} \\
    $0$-$1$ semi-adaptivity (Risky) & \multicolumn{2}{c}{$= 1 + \ln 2 \approx 1.69$} & BTC~\cite{BarakTalgamCohen2026} \\
    \hline
\end{tabular}
\end{center}

\begin{remark}
    Key results from Dean--Goemans--Vondr\'{a}k~\cite{DeanGoemansVondrak2008}:
    \begin{itemize}
        \item Martingale argument (Lemma~3.1): for any adaptive policy, the expected total truncated mass of items attempted is $\le 2$, giving the LP bound $\text{ADAPT} \le \Phi(2)$.
        \item $4$-approximation nonadaptive algorithm using polymatroid LP bound $\Psi(2)$.
        \item Adaptivity gap $\le 4$ for Non-Risky SK.
    \end{itemize}

    Key results from Barak--Talgam-Cohen~\cite{BarakTalgamCohen2026}:
    \begin{itemize}
        \item Introduction of \emph{semi-adaptivity}: a $k$-semi-adaptive policy makes at most $k$ adaptive queries. Provides smooth interpolation between nonadaptive ($k=0$) and fully adaptive ($k=n$).
        \item With $\tilde{O}(1/\eps)$ adaptive queries, achieves a $(6.44 + \sqrt{\eps})$-approximation for Risky SK, improving on the best known fully adaptive $8$-approximation.
        \item The $0$-$1$ semi-adaptivity gap is \emph{exactly} $1 + \ln 2 \approx 1.69$ for Risky SK: even a single observation of remaining capacity gives a provable constant-factor improvement.
    \end{itemize}
\end{remark}

%% ============================================================
\subsubsection{Classical TSP Hardness and Approximation (Background)}
%% ============================================================

For context, we list the key results for deterministic TSP that underlie the stochastic results.

\begin{center}
\renewcommand{\arraystretch}{1.4}
\resizebox{\textwidth}{!}{%
\begin{tabular}{l c l}
    \hline
    \textbf{Result} & \textbf{Bound} & \textbf{Source} \\
    \hline
    TSP is NP-hard & --- & Karp~\cite{Karp1972} \\
    TSP is APX-hard & --- & Karpinski--Lampis--Schmied~\cite{Karpinski2015} \\
    Sym.\ TSP: $3/2$-approx (Christofides) & $3/2$ & Christofides~\cite{Christofides1973} \\
    Sym.\ TSP: $(3/2 - 10^{-36})$-approx & $< 3/2$ & Karlin--Klein--OG~\cite{Karlin2021} \\
    Asym.\ TSP: $O(1)$-approx & $O(1)$ & Svensson--Tarnawski--V\'{e}gh~\cite{Svensson2020} \\
    \hline
\end{tabular}%
}
\end{center}

%% ============================================================
\subsubsection{Summary Table: All Stochastic TSP Gap Bounds}
%% ============================================================

\begin{center}
\renewcommand{\arraystretch}{1.5}
\resizebox{\textwidth}{!}{%
\begin{tabular}{l l c c c}
    \hline
    \textbf{Metric} & \textbf{Gap} & \textbf{Ratio} & \textbf{Best Known Bounds} & \textbf{Source} \\
    \hline
    Symmetric & Obliviousness $\f{\optprior}{\optpost}$ & a priori / a posteriori & $\le 3$ & \cite{Shmoys2008} \\
    Symmetric & Adaptivity $\f{\optprior}{\optadapt}$ & a priori / adaptive & $\le 3$ (inherited) & \cite{Shmoys2008} \\
    Symmetric & Clairvoyance $\f{\optadapt}{\optpost}$ & adaptive / a posteriori & $\le 3$ (inherited) & \cite{Shmoys2008} \\
    Sym.\ tree & All gaps & --- & $= 1$ & \cite{SchalekampShmoys2008} \\[4pt]
    Asymmetric & Obliviousness $\f{\optprior}{\optpost}$ & a priori / a posteriori & $\Omega\!\left(\f{n^{1/4}}{\log n}\right)$, $O(\sqrt{n})$ & \cite{Christalla2025} \\
    Asymmetric & Adaptivity $\f{\optprior}{\optadapt}$ & a priori / adaptive & $\Omega\!\left(\f{n^{1/4}}{\log n}\right)$, $O(\sqrt{n})$ & \cite{Christalla2025} \\
    Asymmetric & Clairvoyance $\f{\optadapt}{\optpost}$ & adaptive / a posteriori & $O(\sqrt{n})$; \textbf{open if $O(1)$} & --- \\
    \hline
\end{tabular}%
}
\end{center}
